\documentclass[11pt]{article}
\usepackage{amsmath,amssymb,amsthm,hyperref}
\newtheorem{theorem}{Theorem}
\newtheorem{corollary}[theorem]{Corollary}
\title{Squarefree Reciprocal Bases and Finite-Capacity Averaging}
\author{}
\date{November 23, 2026}
\begin{document}
\maketitle

\paragraph{Scope.}
These are exact algebraic constructions on the Boolean cube. They
compress a complete reciprocal lattice and give positive cancellation
identities, but they do not prove the requested worst-case head bound.
The proofs here are manual, not Lean formalized.

\section{A basis indexed by small subsets}

Fix $0\le k\le n$, $\delta>0$, and an affine function $C$ that is
positive on the Boolean cube. Put $z_i=1-x_i$ and
\[
 Q(t)=\prod_{j=0}^k(t+\delta j),\qquad
 B_S(x)=C(x)+\delta\left(k-\sum_{i\in S}z_i\right)
          \quad (S\subseteq[n],\ |S|\le k).
 \tag{1}
\]
Every $B_S$ is positive on the cube. If $C$ has strictly positive
variable coefficients, all the $B_S$ satisfy the manuscript's
strictly positive slope condition as well. Let $\mathcal P_d$
denote the space of Boolean polynomial tables of degree at most $d$,
and put $D_d=\sum_{j=0}^{\min(d,n)}\binom nj$.

\begin{theorem}[Exact squarefree basis]
The $D_k$ tables $1/B_S$ in (1) form a basis of
\[
                 Q(C)^{-1}\mathcal P_k.
 \tag{2}
\]
Consequently their affine-numerator span is exactly
$Q(C)^{-1}\mathcal P_{k+1}$. In particular, $1/C$ is a unique
constant-coefficient linear combination of the $1/B_S$.
\end{theorem}
\begin{proof}
For $S$ of size $m$, take its subset difference
\[
 U_S=\sum_{T\subseteq S}(-1)^{m-|T|}/B_T.
 \tag{3}
\]
If some $z_i=0$ with $i\in S$, terms cancel in pairs. If every
$z_i=1$ on $S$, the elementary reciprocal difference formula gives
\[
 U_S=\frac{m!\delta^m z_S}
                {\prod_{j=0}^m(C+\delta(k-j))},\qquad
 z_S=\prod_{i\in S}z_i.
 \tag{4}
\]
For completeness, the formula
$\sum_{j=0}^m(-1)^{m-j}\binom mj/(t-\delta j)
=m!\delta^m/\prod_{j=0}^m(t-\delta j)$ follows by induction:
the next difference is the current expression at $t-\delta$
minus its value at $t$, and the resulting numerator is multiplied
by $(m+1)\delta$. Here every denominator is positive, since
$t=C+\delta k$ and $m\le k$.

Multiplying (4) by $Q(C)$ gives
\[
 Q(C)U_S=m!\delta^m z_S
                 \prod_{j=0}^{k-m-1}(C+\delta j).
 \tag{5}
\]
The empty product is one. Multilinear reduction has degree at most
$k$, and every monomial contains $S$. Its coefficient of $z_S$ is
\[
 m!\delta^m\prod_{j=0}^{k-m-1}
                  (C(x^{(S)})+\delta j)>0,
 \tag{6}
\]
where $x^{(S)}$ is zero on $S$ and one outside $S$. To see this
coefficient directly, set all outside $z_i$ to zero and all inside
$z_i$ to one in the factor following $z_S$. Reduction modulo
$z_i^2-z_i$ gives exactly (6).

Order subsets by increasing size, breaking ties arbitrarily. The
coefficient matrix of (5) in the monomials $z_T$, $|T|\le k$, is
triangular with nonzero diagonal. Thus these tables form a basis
of $\mathcal P_k$. The subset-difference transform (3) is invertible
by ordinary subset inversion. This proves (2).

Multiplication by any affine numerator raises degree by at most
one. Conversely every monomial of degree at most $k+1$ is an
affine function times a monomial of degree at most $k$ (or is
already in that lower-degree space), proving the span equality.
Finally, $Q(C)/C=\prod_{j=1}^k(C+\delta j)$ has degree at most $k$.
\end{proof}

The complete lattice of indices $\alpha\in\mathbb Z_{\ge0}^{n+1}$
with $|\alpha|=k$ has $\binom{n+k}{k}$ nodes. Every reciprocal
from that lattice lies in (2), by division of $Q(t)$ by
$t+\delta L_\alpha(x)$ as in \texttt{polya-jet-filtration.tex}.
Thus (1) preserves its entire constant-numerator span using only
$D_k$ nodes. It does not preserve every particular positive
representation: the unique coefficients of $1/C$ in this basis
can be negative. For example, at $n=3,k=2,\delta=1$ and
$C=1+x_1+2x_2+3x_3$, the coefficient at $S=\varnothing$ is
$-19/140$. This value was also independently checked by exact
rational interpolation at every vertex.

\section{A generalized binomial averaging identity}

The following generating-function identity includes both reinforced
and finite-capacity positive averaging. Write
\[
 B(x)=b_0+\sum_{i=1}^n b_i x_i,\qquad
 T=\sum_{i=0}^n b_i,
 \qquad b_i>0,\quad\delta>0.
 \tag{7}
\]
For real $a$ and integer $r\ge0$, use the generalized binomial
coefficient $\binom ar=a(a-1)\cdots(a-r+1)/r!$.

\begin{theorem}[Generalized binomial reciprocal identity]
For any nonzero real $s$ and any integer $k\ge0$, at every Boolean
input,
\[
 \frac{\binom{T/s}{k}}{B(x)}
 =\sum_{\substack{\alpha\in\mathbb Z_{\ge0}^{n+1}\\|\alpha|=k}}
 \frac{\left(1+\delta\alpha_0/b_0\right)
                  \prod_{i=0}^n\binom{b_i/s}{\alpha_i}}
      {B(x)+\delta\left(\alpha_0+\sum_{i=1}^n\alpha_i x_i\right)}.
 \tag{8}
\]
No division by $\binom{T/s}{k}$ is required in this statement.
Every denominator on the right is admissible with strictly positive
slopes.
\end{theorem}
\begin{proof}
Fix an input, write $B=B(x)$, and use $1/a=\int_0^1t^{a-1}dt$
for positive $a$. For $|u|<1/2$, the generating function of the
right-hand side coefficients is
\[
 \int_0^1 t^{B-1}(1+u t^\delta)^{B/s-1}
       \left(1+(1+\delta/s)u t^\delta\right)
                       (1+u)^{(T-B)/s}\,dt.
 \tag{9}
\]
Indeed the coordinate-zero factor is
\[
 \sum_{r\ge0}\left(1+\frac{\delta r}{b_0}\right)
                  \binom{b_0/s}{r}v^r
   =(1+v)^{b_0/s-1}\left(1+(1+\delta/s)v\right).
\]
Each active coordinate contributes $(1+u t^\delta)^{b_i/s}$;
each inactive coordinate contributes $(1+u)^{b_i/s}$.
The series are absolutely convergent for the displayed range of $u$,
and their integrands are dominated by a constant times $t^{B-1}$,
so coefficient extraction under the integral is justified.

The part of (9) depending on $t$ is $1/B$ times the derivative of
$t^B(1+u t^\delta)^{B/s}$. Its boundary values are zero and
$(1+u)^{B/s}$. Consequently (9) equals $(1+u)^{T/s}/B$.
Extracting the coefficient of $u^k$ proves (8).
\end{proof}

If $s<0$, every nonzero binomial product on the right of (8) has
sign $(-1)^k$, as does the left numerator. Normalizing therefore
gives strictly positive weights. If $s>0$ and every $b_i/s$ is a
positive integer $q_i$, all weights are nonnegative and supported
on $0\le\alpha_i\le q_i$. For $0\le k\le\sum_i q_i$, the
normalizing coefficient is positive. The number of retained
indices is the coefficient of $u^k$ in
$\prod_{i=0}^n(1+u+\cdots+u^{q_i})$.

\paragraph{An independent finite-population proof.}
For integer capacities $q_i=b_i/s$, give each of $q_i$ tokens color $i$.
Choose a uniform $k$-subset of the $M=\sum_i q_i$ tokens and let
$\alpha_i$ count selected tokens of color $i$. The probability of
$\alpha$ is $\prod_i\binom{q_i}{\alpha_i}/\binom Mk$.
For a fixed input, call colors $0$ and $i$ with $x_i=1$ active.
Conditional on the number $r$ of selected active tokens, the expected
number of selected color-zero tokens is $r b_0/B$. Therefore
\[
 \mathbb E\left[\left.
      \frac{b_0+\delta\alpha_0}{B+\delta r}\right|r\right]
      =\frac{b_0+\delta r b_0/B}{B+\delta r}
      =\frac{b_0}{B}.
\]
Taking expectations gives (8) after normalization. This proves the
finite-capacity case without using power series or integration.
The symmetry underlying this averaging is separately kernel checked
by \texttt{HeadComplexity.permutation\_reciprocal\_ratio\_sum}
in \texttt{Results/PermutationReciprocalAverage.lean}. The conversion
to the binomial count formula above remains a manual argument.

\begin{corollary}[Two adjacent subset levels]
Let $b>0$, $B=b(1+|x|)$, $1\le k\le n$, and
$B_S=B+\delta(k-|S|+\sum_{i\in S}x_i)$. Then
\[
 \frac{\binom{n+1}{k}}{B}
   =\sum_{|S|=k}\frac1{B_S}
       +(1+\delta/b)\sum_{|S|=k-1}\frac1{B_S}.
 \tag{10}
\]
This is a positive identity with exactly $\binom{n+1}{k}$ terminal
forms. All denominators have strictly positive variable coefficients.
\end{corollary}
\begin{proof}
In (8) set $b_i=s=b$ for every $i$. Then each $\alpha_i$ is zero
or one; $\alpha_0=0$ gives subsets of size $k$, and
$\alpha_0=1$ gives subsets of size $k-1$.
\end{proof}

At $k=n+1$ the same identity has a single terminal denominator
proportional to $B$, so it is merely a rescaling. That endpoint
has no encoding consequence. More generally, (8) is an exact
cancellation identity, not an independently programmable sign gadget.
The common-lattice filtration and its restrictions on low-order
variations still apply. No compatibility theorem converting these
identities into all Boolean sign tables at the desired head count
is supplied here.

\end{document}
